% Źródła: Eves s. 288, 291--295; Wikipedia: Hyperbolic geometry, Non-Euclidean geometry.

\section{Równoległe i hiperrównoległe} % lang-pl
\section{Parallele e iperparallele} % lang-it
\label{sekcja_hiperboliczna}

Zobaczymy teraz, co odkrywcy zobaczą pierwsze, gdy zastąpią piąty postulat jego zaprzeczeniem. % lang-pl
W tej sekcji wskażemy, jak wygląda przez taki punkt zbiór prostych równoległych i jak zachowuje się równoległość. % lang-pl
Vedremo ora ciò che gli scopritori vedranno per primo quando sostituiranno il quinto postulato con la sua negazione. % lang-it
In questa sezione indicheremo come si presenta, per un tale punto, l'insieme delle rette parallele e come si comporta il parallelismo. % lang-it

\begin{definition}
[geometria hiperboliczna] % lang-pl
[geometria iperbolica] % lang-it
	Geometria płaska oparta na aksjomatach geometrii neutralnej i na zaprzeczeniu aksjomatu Playfaira: przez punkt leżący poza prostą przechodzi więcej niż jedna prosta do niej równoległa. % lang-pl
	Geometria piana fondata sugli assiomi della geometria neutrale e sulla negazione dell'assioma di Playfair: per un punto esterno a una retta passa più di una retta a essa parallela. % lang-it
\index{geometria!hiperboliczna} % lang-pl
\index{geometria!iperbolica} % lang-it
\end{definition}

Nazwę \emph{hiperboliczna} wprowadzi Felix Klein w 1871 roku (zobacz sekcję \ref{sekcja_beltrami_cayley_klein}). % lang-pl
O geometrii hiperbolicznej wspomina się w $\Delta_{24}^7$. % lang-pl
Il nome \emph{iperbolica} lo introdurrà Felix Klein nel 1871 (si veda la sezione \ref{sekcja_beltrami_cayley_klein}). % lang-it
La geometria iperbolica è menzionata in $\Delta_{24}^7$. % lang-it

Zakładamy, że wszystko, co wynika z aksjomatów bez piątego postulatu -- w tym pierwsze dwadzieścia osiem twierdzeń pierwszej księgi Euklidesa -- pozostaje w mocy \cite[s. 288, 291]{eves1_1972}. % lang-pl
Do tego dodajemy jeden aksjomat (Eves przyjmuje go w postaci mocniejszej, niż jest potrzebna, dla wygody): % lang-pl
Assumiamo che tutto ciò che discende dagli assiomi senza il quinto postulato -- comprese le prime ventotto proposizioni del primo libro di Euclide -- resti valido \cite[p. 288, 291]{eves1_1972}. % lang-it
A ciò aggiungiamo un solo assioma (Eves lo assume in una forma più forte del necessario, per comodità): % lang-it

\begin{axiom}
[Łobaczewskiego] % lang-pl
[di Lobachevskij] % lang-it
\index{aksjomat!Łobaczewskiego} % lang-pl
\index{assioma!di Lobachevskij} % lang-it
	Niech punkt $P$ nie leży na prostej $AB$, a $Q$ będzie spodkiem prostopadłej z $P$ na $AB$. % lang-pl
	Istnieją wtedy dwie półproste $PX$ i $PY$, niezawarte w jednej prostej i nieprzecinające $AB$, takie że każda półprosta z $P$ leżąca wewnątrz kąta $XPY$ zawierającego $PQ$ przecina $AB$. % lang-pl
	Sia $P$ un punto non appartenente alla retta $AB$ e sia $Q$ il piede della perpendicolare da $P$ a $AB$. % lang-it
	Esistono allora due semirette $PX$ e $PY$, non contenute in una stessa retta e che non intersecano $AB$, tali che ogni semiretta uscente da $P$ e interna all'angolo $XPY$ contenente $PQ$ interseca $AB$. % lang-it
\end{axiom}

Eves \cite[s. 291]{eves1_1972}. % lang-pl
Eves \cite[p. 291]{eves1_1972}. % lang-it

Półproste $PX$ i $PY$ nazwiemy \emph{równoległymi} do $AB$ przez $P$, a proste przez $P$ leżące poza kątem $XPY$ zawierającym $PQ$ -- \emph{hiperrównoległymi} do $AB$ \cite[s. 291--292]{eves1_1972}. % lang-pl
Le semirette $PX$ e $PY$ le chiameremo \emph{parallele} ad $AB$ per $P$, e le rette per $P$ esterne all'angolo $XPY$ contenente $PQ$ \emph{iperparallele} ad $AB$ \cite[p. 291--292]{eves1_1972}. % lang-it
\index{proste!hiperrównoległe} % lang-pl
\index{rette!iperparallele} % lang-it

\begin{theorem}
	Kąty $XPQ$ i $YPQ$ są równe i ostre. % lang-pl
	Gli angoli $XPQ$ e $YPQ$ sono uguali e acuti. % lang-it
\end{theorem}

Dowód: Eves \cite[s. 292]{eves1_1972}. % lang-pl
Dimostrazione: Eves \cite[p. 292]{eves1_1972}. % lang-it

Kąt $XPQ$ nazwiemy \emph{kątem równoległości}; wynika z niego też, że przez $P$ przechodzi nieskończenie wiele prostych hiperrównoległych do $AB$ \cite[s. 292]{eves1_1972}. % lang-pl
L'angolo $XPQ$ si chiamerà \emph{angolo di parallelismo}; ne segue anche che per $P$ passano infinite rette iperparallele ad $AB$ \cite[p. 292]{eves1_1972}. % lang-it
\index{kąt!równoległości} % lang-pl
\index{angolo!di parallelismo} % lang-it

\begin{proposition}
	Równoległość skierowanych prostych jest przenoszalna (jeśli $\overline{PX} \parallel \overline{AB}$, to $\overline{RX} \parallel \overline{AB}$ dla każdego $R$ na $PX$ po stronie $X$), symetryczna i przechodnia. % lang-pl
	Il parallelismo tra rette orientate è trasmissibile (se $\overline{PX} \parallel \overline{AB}$, allora $\overline{RX} \parallel \overline{AB}$ per ogni $R$ su $PX$ dalla parte di $X$), simmetrico e transitivo. % lang-it
\end{proposition}

Dowody: Eves \cite[s. 292--295]{eves1_1972}; dowód symetrii Eves przypisuje Łobaczewskiemu, a dowód przechodniości (gdy dwie proste leżą po jednej stronie trzeciej) Gaussowi. % lang-pl
Dimostrazioni: Eves \cite[p. 292--295]{eves1_1972}; la dimostrazione della simmetria Eves l'attribuisce a Lobachevskij, quella della transitività (quando due rette stanno dalla stessa parte della terza) a Gauss. % lang-it

Zadania do tej sekcji: Eves \cite[s. 288--291]{eves1_1972}. % lang-pl
Esercizi per questa sezione: Eves \cite[p. 288--291]{eves1_1972}. % lang-it
